% Template for post or presentation abstracts submitted to ILAS 2022
% 1. Fill in the presenter's information (name, affiliation, email
% address) and talk information (title of presentation, abstract).
% 2. Compile to make sure there are no errors.
% 3. Save the .tex file for your abstract with a distinctive name,
% ideally "FamilynameGivenname.tex".
% 4. Upload your .tex file to https://clr.ie/129557
%       
% * Please keep your LaTeX commands simple, and avoid the use of
%    user-defined macros.
% * Include details of co-authors and funding acknowledgements after the abstract.
% * Upload only the LaTeX file (with .tex extension).
% * Questions: Contact Niall Madden (Niall.Madden@NUIGalway.ie)

\documentclass[a4paper]{amsart} 
\usepackage{amssymb} 
\usepackage{hyperref}
\pagestyle{empty}
\date{} 

%% IGNORE THE NEXT 4 LINES
\makeatletter % 
\def\labelenumi{\theenumi.}
\def\theenumi{\@arabic\c@enumi}
\makeatother
%% STOP IGNORING NOW

\begin{document}

\title{Monge-like properties in the interval setting}
\author{\v{C}ern\'{y} Martin} %Presenting author only
\address{Charles University} % Affiliation of presenting author
\email{cerny@kam.mff.cuni.cz} % Email address of presenting author only

\maketitle

% The abstract  ***no more than one page***

The Monge property of a real matrix $A \in \mathbb{R}^{n \times n}$ can be expressed as
\[
a_{ij} + a_{k\ell} \leq a_{i\ell} + a_{kj} 
\]
for $1 \leq i \leq m$, $1 \leq j \leq n$. This simple to express property is fundamental for many efficient algorithms applied in various problems from geometry, combinatorics, optimisation or statistics~\cite{survey,survey2}. Many modifications of the property lead to interesting classes of Monge-like matrices, to name a few, \textit{Robinsonian}, \textit{ultrametric}, \textit{totally monotonic}, or \textit{Monge-permutable} matrices.

In our research, we do a systematic analysis of Monge-like properties in the interval setting. Matrix intervals allow for computations with inexact data. Rather than storing precise values (which might be impossible thanks to limitations of the measurement), we employ real intervals with a guarantee that each underlying value is in the range of its interval. Computations with matrix intervals are then carried out in a way that the exact solution of the original problem is guaranteed to be in the range of the output interval. Further, if the width of the interval is negligible, so is the error. 

We focus on two variants of the Monge-like properties in the interval setting, so called \textit{weak} and \textit{strong} properties. For the definition of both of these properties, matrix realisations of the matrix interval (real matrices with entries from the intervals) are considered. If there is at least one realisation satisfying the Monge-like property, we say the matrix interval satisfies its weak form. If all of the realisations satisfy the Monge-like property, the matrix interval satisfies its strong form.

We deal with different characterisations of weak and strong Monge-like properties as well as with necessary and sufficient conditions. If the matrix interval satisfies the strong property, an interesting question to consider is if it can be recognised by checking the Monge-like property for a finite (hopefully, polynomial) number of its matrix realisations. This property of matrix intervals, referred to as the \textit{interval property}~\cite{Garloff}, is also studied in our research. Finally, we investigate possible interval generalisations of known algorithm for Monge-like properties together with their complexity analysis.

%% Coauthor details here (optional - delete if not applicable)
% and funding acknowledgment (optional - delete if not applicable)
\emph{This is joint work with prof. J\"{u}rgen Garloff (Konstanz).
Supported by SVV 260578/2020.}

\begin{thebibliography}{1}
\bibitem{survey}
Rainer E. Burkard and Bettina Klinz and R\"{u}diger Rudolf.
\newblock Perspectives of Monge properties in optimization
\newblock {\em  Discrete Applied Mathmematics}, 70:95--161, 1996.

\bibitem{survey2}
Swetha Sethumadhavan.
\newblock A survey of Monge properties
\newblock {\em  UNLV Theses, Dissertations,	Professional Papers, and Capstones}, 1198, 2009.

\bibitem{Robinson}
W. S. Robinson.
\newblock A Method for Chronologically Ordering Archaeological Deposits
\newblock {\em  American Antiquity}, 16(4):293--301, 1951.


\bibitem{Cerny}
Martin \v{C}ern\'{y}.
\newblock Interval matrices with Monge properties
\newblock {\em  Applications of Mathematics}, 65(5):619--643, 2020.

\bibitem{Garloff}
J\"{u}rgen Garloff and Doaa Al-Saafin and Mohammad Adm.
\newblock Further Matrix Classes Posessing Interval Property
\newblock {\em  Konstanzer Schriften in Mathematik}, 398, 2021.


\end{thebibliography}


\end{document}


% Please leave the next bit alone  
%%% Local Variables: 
%%% mode: latex
%%% TeX-master: "../ILAS-2022"
%%% End:
